%%%%%% Verander de volgende lijnen niet %%%%%%%

\documentclass[12pt,a4paper,svgnames,oneside]{book}


\newcommand{\initialiseer}[2]{\newwrite\alleoplossingen
\newwrite\oplossingenzonderster
\immediate\openout\alleoplossingen=#1
\immediate\openout\oplossingenzonderster=#2}

%\immediate\write\oplossingenzonderster{\unexpanded{\renewcommand{\thesection}{}}}

\newcommand{\oplos}[1]{\immediate\write\alleoplossingen{\unexpanded{\begin{opgelost}}{\theoefen} \unexpanded{#1 \end{opgelost}}}\immediate\write\oplossingenzonderster{\unexpanded{\begin{opgelost}}{\theoefen} \unexpanded{#1 \end{opgelost}}}}

\newcommand{\steroplos}[1]{\immediate\write\alleoplossingen{\unexpanded{\begin{opgelostster}}{\theoefen} \unexpanded{#1 \end{opgelostster}}}}

\newcommand{\sluitfiles}{\immediate\closeout\alleoplossingen
\immediate\closeout\oplossingenzonderster} 
\initialiseer{inputs/oplossingen-alle.tmp}{inputs/oplossingen-zonder-ster.tmp}
\usepackage{etex}
\usepackage[bookmarks=false]{hyperref}
\usepackage{easybmat}
\usepackage{fancyhdr}
\usepackage[dutch]{babel}
\usepackage[explicit]{titlesec}
\usepackage{pgf,tikz}
\usepackage{amsmath,amssymb,amsthm}
\usepackage{graphicx}
\usepackage{ifthen}
\usepackage{enumitem}
\usepackage{calc}
\usepackage{eurosym}
\usepackage{pdfpages}

\setlength{\oddsidemargin}{-15pt}
\setlength{\evensidemargin}{0pt}
\setlength{\topmargin}{-50pt}
\setlength{\textheight}{680pt}
\setlength{\textwidth}{485pt}
\setlength{\headsep}{30pt}
\setlength{\parindent}{0pt}
\setlength{\parskip}{1ex plus 0.5ex minus 0.2ex}
\setlength{\footskip}{40pt}
\setlength{\headheight}{50pt}
\addtolength{\baselineskip}{20pt}

\numberwithin{equation}{chapter}
\renewcommand{\theequation}{\arabic{equation}}

\newcommand{\bfit}{\bf\fontshape{it}\selectfont}

%%%%%%%%% De definities voor de headers %%%%%%%%%%%%%%%%%%%%%
%%%%%%%%% Juiste definitie voor input logo geven %%%%%%%%%%%%%%%%%%%%%%%%

\newcommand{\titel}{Je hartslag monitor altijd op zak}

\newcommand{\filenaamlogo}{inputs/logo-innovationlab.png}
%\parbox[position][height][inner-pos]{width}{text}
\newcommand{\insertlogo}{\parbox[c][0.95cm][c]{0.85cm}{\includegraphics[width=0.85cm]{\filenaamlogo}}}

\definecolor{fac-color}{cmyk}{0.65,0,0,0.15}
\definecolor{fac-color-dark}{cmyk}{0.65,0,0,0.35}
\definecolor{fac-color-orange}{cmyk}{0,0.38,0.99,0.13}

\newlength{\evenlengte}

\newcommand*\runningheader[1]{\settowidth{\evenlengte}{\mbox{\bf #1}}\addtolength{\evenlengte}{6mm}%
  \begin{tikzpicture}
    \fill[fill=fac-color] (0,0) rectangle (\textwidth,1.3cm);
    \fill[fac-color-dark] (\textwidth-\evenlengte,0) rectangle (\textwidth,1.3cm);
    \node[xshift=1mm,anchor=south west]{\insertlogo};
    \node[xshift=1.5cm,yshift=5.5mm,anchor=south west]{\color{white}\fontsize{11pt}{13pt}\selectfont\bf Innovation Lab KU Leuven};
    \node[xshift=1.5cm,yshift=0.5mm,anchor=south west]{\color{white}\fontsize{11pt}{13pt}\selectfont\bf \titel};
    \node[xshift=\textwidth-2mm,yshift=0.5mm,anchor=south east]{{\color{white}\fontsize{11pt}{13pt}\selectfont\bf\boldmath #1}};
    \node[xshift=\textwidth-2mm,yshift=6mm,anchor=south east]{\color{white}\fontsize{11pt}{13pt}\selectfont\bf \thepage};
  \end{tikzpicture}}

\newcommand*\hoofdstukheader{%
  \begin{tikzpicture}
    \fill[fill=fac-color] (0,0) rectangle (\textwidth,1.3cm);
    \node[xshift=1mm,anchor=south west]{\insertlogo};
    \node[xshift=1.5cm,yshift=5.5mm,anchor=south west]{\color{white}\fontsize{11pt}{13pt}\selectfont\bf Innovation Lab KU Leuven};
    \node[xshift=1.5cm,yshift=0.5mm,anchor=south west]{\color{white}\fontsize{11pt}{13pt}\selectfont\bf \titel};
    \node[xshift=\textwidth-2mm,yshift=6mm,anchor=south east]{\color{white}\fontsize{11pt}{13pt}\selectfont\bf \thepage};
  \end{tikzpicture}}

\pagestyle{fancy}
\lfoot{}
\cfoot{}
\rfoot{}
\lhead{\runningheader{\nouppercase{\rightmark}}}
\chead{}
\rhead{}

\fancypagestyle{plain}{
   \fancyhf{}
   \lhead{\hoofdstukheader}
   \lfoot{}
   \cfoot{}
   \rfoot{}}



\renewcommand{\headrulewidth}{0pt}
\renewcommand{\footrulewidth}{0pt}

%%%%%%%%%% De definities voor de hoofdstukpagina's %%%%%%%%%%%%%%%%%%%%%%%%%%%%%%%%%

\newcommand*\chapterlabel{}

\newcommand{\chapterskip}{7mm}

\titlespacing*{\chapter}{0pt}{0pt}{0pt}

\titleformat{\chapter}[block]
  {\gdef\chapterlabel{} \huge\bfseries\boldmath}
  {\gdef\chapterlabel{\chaptertitlename}}{0pt}
  {\ifthenelse{\equal{\chapterlabel}{}}{\hoofdstukzondernaam{#1}}{\hoofdstukmetnaam{\chapterlabel}{#1}}
  \vspace{\chapterskip}}

\newcommand{\mijnchaptermark}{}
\renewcommand{\chaptermark}[1]{\markright{\mijnchaptermark}}

\newcommand{\hoofdstukzondernaam}[1]{%
\mbox{}\hspace{3.7cm}
\raisebox{20mm}[0cm]{\begin{tikzpicture}[baseline=(current bounding box.north)]
 \node[anchor=north west,inner sep=4mm,very thick,draw=fac-color-dark,text badly ragged,text width=12cm]{#1};
\end{tikzpicture}}%\vspace{1mm}
}

\newcommand{\hoofdstukmetnaam}[2]{%
\mbox{}\hspace{3.7cm}
\raisebox{20mm}[0cm]{\begin{tikzpicture}[baseline=(current bounding box.north)]
 \node[anchor=south west,rectangle,fill=fac-color-dark]{\color{white}#1};
 \node[anchor=north west,inner sep=4mm,very thick,draw=fac-color-dark,text badly ragged,text width=12cm]{#2};
\end{tikzpicture}}%\vspace{1mm}
}


%%%%%%%%%%% De definities voor de stellingen, etc. %%%%%%%%%%%%%%%%%%%%%%%%%%%%%%%%%%%%%%

\newcommand{\stellingnaam}[1]{\begin{tikzpicture}[baseline=(X.base)] \node[inner sep=0.6mm,rectangle,fill=fac-color-dark](X){\color{white}#1}; \end{tikzpicture}}

\newtheoremstyle{mijnstelling}% name
  {2ex}%      Space above, empty = `usual value'
  {2ex}%      Space below
  {\itshape}% Body font
  {}%         Indent amount (empty = no indent, \parindent = para indent)
  {\bfseries}% Thm head font
  {}%        Punctuation after thm head
  {2ex}%     Space after thm head: " " = normal interword space;
        %       \newline = linebreak
  {\ifthenelse{\equal{#2}{}}{\stellingnaam{#1}}{\stellingnaam{#1 #2}}}% Thm head spec

\newtheoremstyle{mijndefinitie}% name
  {2ex}%      Space above, empty = `usual value'
  {2ex}%      Space below
  {\rm}% Body font
  {}%         Indent amount (empty = no indent, \parindent = para indent)
  {\bfseries}% Thm head font
  {}%        Punctuation after thm head
  {2ex}%     Space after thm head: " " = normal interword space;
        %       \newline = linebreak
  {\ifthenelse{\equal{#2}{}}{\stellingnaam{#1}}{\stellingnaam{#1 #2}}}% Thm head spec

\newtheoremstyle{mijnvoorbeeld}% name
  {2ex}%      Space above, empty = `usual value'
  {2ex}%      Space below
  {\rm}% Body font
  {}%         Indent amount (empty = no indent, \parindent = para indent)
  {\it}% Thm head font
  {.}%        Punctuation after thm head
  {2ex}%     Space after thm head: " " = normal interword space;
        %       \newline = linebreak
  {}% Thm head spec

{\theoremstyle{mijndefinitie}
\newtheorem*{definitie}{Definitie}}

{\theoremstyle{mijnopmerking}
	\newtheorem*{opmerking}{Opmerking}}

{\theoremstyle{mijnopmerking}
	\newtheorem*{notatie}{Notatie}}

{\theoremstyle{mijnopmerking}
	\newtheorem*{onthoud}{Onthoud}}

\newenvironment{ftonthoud}
{\begin{mdframed}\begin{onthoud}}
		{\end{onthoud}\end{mdframed}}

{\theoremstyle{mijnstelling}
\newtheorem{eigenschap}{Eigenschap}[chapter]}
\renewcommand{\theeigenschap}{\arabic{eigenschap}}

{\theoremstyle{mijnvoorbeeld}
	\newtheorem{voorbeeld}{Voorbeeld}[subsection]}
\renewcommand{\thevoorbeeld}{\arabic{voorbeeld}}

\newenvironment{bewijs eig}{\begin{proof}[Bewijs]}{\end{proof}}
\newenvironment{bewijs oplossing}{\begin{proof}[Uitgeschreven bewijs]}{\end{proof}}

%%%%%%%%%%%%% Abstracts van de hoofdstukken %%%%%%%%%%%%%%%%%%%%%%%%%%%%%%%

\newlength{\breedtesamenvatting}
\setlength{\breedtesamenvatting}{430pt}
\newlength{\samenvattingmarge}
\setlength{\samenvattingmarge}{(\textwidth - \breedtesamenvatting)/2}

\newenvironment{samenvatting}{\par\mbox{}\hspace{\samenvattingmarge}\begin{minipage}[t]{\breedtesamenvatting}}{\end{minipage}\vspace{3ex}\par}


%%%%%%%%%%%%% Er zijn geen hoofdstukken, maar wel module's, intermezzo's, etc, vandaar deze definities %%%%%%%%%%%%%%%%%%%%%%%%%%
%%%%%%%%%%%%% Het \deelmetoef veronderstelt dat de oplossingenfile's in gebruik zijn %%%%%%%%%%%%%%%%%%%%%%%%%%%%%%%%%%%%%%%%%%%%
%%%% De structuur is respectievelijk %%%
%%%%
%%%% \deelmetoef[optioneel argument met korte running header]{benaming van het deel met het nummer}{titel van het deel}{titel in table of contents}{titel in de oplossingenlijst}{running header in oplossingenfile} %%%
%%%%
%%%% \deelzonderoef[optioneel argument met korte running header]{benaming van het deel met het nummer}{titel van het deel}{titel in table of contents} %%%%
%%%%
%%%% \voorbereidendesessie[optioneel argument met korte running header]{benaming van de voorbereidende sessie}
%%%%
%%%% \oplossingenvoorbereidendesessie{Oplossingen van de ...}
%%%%

\renewcommand{\thechapter}{}

\newcommand{\mijnoplosmark}{}

\newcommand{\deelmetoef}[6][xyza]{\renewcommand{\chaptertitlename}{#2}%
\ifthenelse{\equal{#1}{xyza}}{\renewcommand{\mijnchaptermark}{#4}}{\renewcommand{\mijnchaptermark}{#1}}%
\chapter[#4]{#3}%
\immediate\write\alleoplossingen{\unexpanded{\renewcommand{\mijnoplosmark}}{#6}\unexpanded{\section}[#4]{#5}\unexpanded{\markright}{#6}}%
\immediate\write\oplossingenzonderster{\unexpanded{\renewcommand{\mijnoplosmark}}{#6}\unexpanded{\section}[#4]{#5}\unexpanded{\markright}{#6}}
}

\newcommand{\deelzonderoef}[4][xyza]{\renewcommand{\chaptertitlename}{#2}%
\ifthenelse{\equal{#1}{xyza}}{\renewcommand{\mijnchaptermark}{#4}}{\renewcommand{\mijnchaptermark}{#1}}%
\chapter[#4]{#3}}


\newcommand{\voorbereidendesessie}[2][xyza]{\renewcommand{\chaptertitlename}{Voorbereidende sessie}%
\ifthenelse{\equal{#1}{xyza}}{\renewcommand{\mijnchaptermark}{Voorbereidende sessie. #2}}{\renewcommand{\mijnchaptermark}{Voorbereidende sessie. #1}}%
\chapter[#1]{#2}
}

\newcommand{\oplossingenvoorbereidendesessie}[2]{\renewcommand{\chaptertitlename}{}\chapter{#1}\markright{#2}}

%%%%%%%%%%%% Figuren met een ondertitel %%%%%%%%%%%%%%%%%%%%%%%%%%%%%%%%%%%%%%%%%%%%%%%%%%%%%%%%%%%%%%%%%%%%%%%%%%%%%%%%%%%%%%%%%%%%%%%%%%%%%%%
%%%%%% De structuur is als volgt: \figuurmetlabel[\label{ref-figuur} - optioneel]{width=7cm}{filenaam.pdf}{Op deze figuur zie je...} %%%%%%%%%%
%%%%%% Gewone figuur: \gewonefiguur{width=7cm}{filenaam.pdf} %%%%%%%%%%%%%%%%%%%%%%%%%%%%%%%%%%%%%%%%%%%%%%%%%%%%%%%%%%%%%%%%%%%%%%%%%%%%%%%%%%

\newcounter{figuur}[chapter]
\renewcommand{\thefiguur}{\arabic{figuur}}
\newcommand{\figuurmetlabel}[4][]{\ifthenelse{\equal{#4}{}}%
	{\begin{center}\includegraphics[#2]{inputs/#3}\nopagebreak\par\refstepcounter{figuur}Figuur \thefiguur#1\end{center}}%
	{\begin{center}\includegraphics[#2]{inputs/#3}\nopagebreak\par\refstepcounter{figuur}Figuur \thefiguur.\ #4#1\end{center}}%
}


\newcommand{\gewonefiguur}[2]{\begin{center}\includegraphics[#1]{inputs/#2}\end{center}}

%%%%%%%%%%%% Tabellen %%%%%%%%%%%%%%%%%%%%%%%%%%%%%%%%%%%%%%%%%%%%%%%%%%%%%%%%%%%%%%%%%%%%%%%%%%%%%%%%%%%%%%%%%%%%%%%%%%%%%%%%%%%%%%%%%%%%%%%%%%%%%%%%%%
%%%%%% De structuur is als volgt: \begin{tabel}{Al of niet de ondertitel van de tabel}\begin{tabular}[pos]{table spec} ... \end{tabular}\end{tabel} %%%%

\newcounter{tabelnr}[chapter]
\newcommand{\tabeltitel}{}
\renewcommand{\thetabelnr}{\arabic{tabelnr}}
\newenvironment{tabel}[1]{\begin{center}\refstepcounter{tabelnr}\renewcommand{\tabeltitel}{#1}}{\ifthenelse{\equal{\tabeltitel}{}}%
{\nopagebreak\vspace{1ex}\par Tabel \thetabelnr\end{center}}%
{\nopagebreak\vspace{1ex}\par Tabel \thetabelnr.\ \tabeltitel\end{center}}}

%%%%%%%%% Instelling van de opsommingen %%%%%%%%%%%%%%%

\setenumerate{leftmargin=5ex,labelsep=*,label=\arabic*.}

%%%%%%%%% Omgevingen waarin de eningen (oef en steroef) worden getypeset. %%%%%%%
%%%%%%%%% De omgevingen opgelost en opgelostster worden ingeroepen in de automatisch gegenereerde oplossingenfile om de oplossingen te typesetten. %%%%%%%%%%%

\newcommand{\kadertje}[1]{\begin{tikzpicture}[baseline=(X.base)] \node[inner sep=1mm,rectangle,fill=fac-color-dark](X){\mbox{\color{white}\bf#1}}; \end{tikzpicture}}

\newlength{\oploslabel}
\setlength{\oploslabel}{8mm}
\newlength{\labelafstand}
\setlength{\labelafstand}{3mm}

\newenvironment{opgelost}[1]{\begin{list}{}{\setlength{\labelwidth}{\oploslabel}\setlength{\leftmargin}{\oploslabel}\addtolength{\leftmargin}{\labelafstand}\setlength{\labelsep}{\labelafstand}} \item[{\parbox[t]{\oploslabel}{\kadertje{#1}}}] }{\end{list}}

\newenvironment{opgelostster}[1]{\begin{list}{}{\setlength{\labelwidth}{\oploslabel}\setlength{\leftmargin}{\oploslabel}\addtolength{\leftmargin}{\labelafstand}\settowidth{\labelsep}{{\boldmath$\ast$}}\addtolength{\labelsep}{\labelafstand}} \item[{\parbox[t]{\oploslabel}{\mbox{}{\color{fac-color-dark}{\boldmath$\ast$}}\kadertje{#1}}}] }{\end{list}}

\newcounter{oefen}[chapter]
\renewcommand{\theoefen}{\arabic{oefen}}

\newenvironment{somop}{\begin{list}{\parbox[t]{\oploslabel}{$\bullet$}}{\setlength{\labelwidth}{\oploslabel}\setlength{\leftmargin}{\oploslabel}\addtolength{\leftmargin}{\labelafstand}\setlength{\labelsep}{\labelafstand}}}{\end{list}}

\newenvironment{oef}{\begin{list}{}{\setlength{\labelwidth}{\oploslabel}\setlength{\labelsep}{\labelafstand}\setlength{\leftmargin}{\oploslabel}\addtolength{\leftmargin}{\labelafstand}}\refstepcounter{oefen} \item[{\parbox[t]{\oploslabel}{\kadertje{\theoefen}}}] }{\end{list}}

\newenvironment{steroef}{\begin{list}{}{\setlength{\labelwidth}{\oploslabel}\setlength{\leftmargin}{\oploslabel}\addtolength{\leftmargin}{\labelafstand}\settowidth{\labelsep}{{\boldmath$\ast$}}\addtolength{\labelsep}{\labelafstand}}\refstepcounter{oefen} \item[{\parbox[t]{\oploslabel}{\mbox{}{\color{fac-color-dark}{\boldmath$\ast$}}\kadertje{\theoefen}}}] }{\end{list}}

%%%%%%%%%%% Sectietitels. Het commando \zondernummer wordt aangeroepen bij de opmaak van de %%%%%%%%%%%%%
%%%%%%%%%%% sectietitels bij de oplossingen. %%%%%%%%%%%%%%%%%%%%

\renewcommand{\thesection}{\arabic{section}}

\renewcommand{\sectionmark}[1]{\markright{#1}}

\newlength{\sectielabelbreedte}
\newlength{\sectielabelskip}
\newlength{\beschikbaar}
\newlength{\streepdikte}

\setlength{\sectielabelbreedte}{6mm}
\setlength{\sectielabelskip}{4mm}
\setlength{\streepdikte}{2mm}

\titleformat{\section}[block]{}{}{0pt}{\metnummer{\thesection}{#1}}

\newcommand{\metnummer}[2]{\setlength{\beschikbaar}{\textwidth-\sectielabelbreedte-\sectielabelskip-4mm}%
\begin{tikzpicture}
\fill[fill=fac-color-dark] (0,0) rectangle (\textwidth,\streepdikte);
\node[anchor=north west,fill=fac-color-dark,rectangle,inner sep=2mm]{\parbox[t]{\sectielabelbreedte}{\Large\bf\color{white}#1}};
\node[anchor=north east,xshift=\textwidth,yshift=-\streepdikte,inner sep=0pt]{\parbox[t]{\beschikbaar}{\Large\bf\boldmath#2}};
\end{tikzpicture}}

\newlength{\streeplengte}
\setlength{\streeplengte}{9mm}

\newcommand{\zondernummer}[1]{\setlength{\beschikbaar}{\textwidth-2\streepdikte}%
\settoheight{\streeplengte}{\parbox[b]{\beschikbaar}{\Large\bf#1}}%
\addtolength{\streeplengte}{4mm}%
\begin{tikzpicture}
\fill[fill=fac-color-dark] (0,0) rectangle (\textwidth,\streepdikte);
\fill[fill=fac-color-dark] (0,0) rectangle (\streepdikte,-\streeplengte);
\node[anchor=north east,xshift=\textwidth,yshift=-\streepdikte,inner sep=0pt]{\parbox[t]{\beschikbaar}{\Large\bf\boldmath#1}};
\end{tikzpicture}}


%%%%%%%% Formattering subsections %%%%%%%%%%%

%\titleformat{\subsection}[block]{}{}{0pt}{\large\bf\boldmath#1}

\titleformat{\subsection}[block]{\bf\boldmath}{\thesubsection}{1.5ex}{#1}

%%%%%%%%%%%%%%%%%%%%%%%%%%%%%%%%%%%%%%%%%%%%%%%%%%
%%% Een specialleke bij Google: in te leveren opdracht %%%
%%%%%%%%%%%%%%%

\newenvironment{opdracht}[1]{\mijnopdracht{#1}}{\vspace{3ex}\par}

\newcommand{\mijnopdracht}[1]{{\titleformat{\subsection}[block]{}{}{0pt}{\opdrachttypeset{#1}}%
\subsection*{#1}}}

\newlength{\overschot}
\newlength{\linkerlijn}
\newlength{\eenbreedte}
\setlength{\linkerlijn}{1cm}

\newcommand{\opdrachttypeset}[1]{\settowidth{\eenbreedte}{\bf\large\boldmath#1}%
\setlength{\overschot}{\textwidth-\linkerlijn-\eenbreedte-4ex}%
\begin{tikzpicture}
\fill[fill=fac-color-orange] (0,\streepdikte/2) rectangle (\linkerlijn,-\streepdikte/2);
\node at (\linkerlijn+2ex,0) [anchor=west,inner sep=0pt] {{\bf\large\boldmath#1}};
\fill[fill=fac-color-orange] (\textwidth-\overschot,\streepdikte/2) rectangle (\textwidth,-\streepdikte/2);
\end{tikzpicture}}

%%%%%%%%%%%%%%%%%%%%%%%%%%%%%%%%%%%%%%%%%%%%%%%%%%%%%

\usepackage{amsmath}
\usepackage{hyperref}
\usepackage{subcaption}
\usepackage{tikz}
\usepackage{pgfplots,mathrsfs}
\usepackage{marginnote}
\usepackage[bottom]{footmisc}
\usepackage{schemabloc}
\usetikzlibrary{circuits}
\usepackage{verbatim}
\usetikzlibrary{arrows, decorations.markings}

\begin{document}
\addtocounter{page}{-1}

%% We beginnen met het logo te defini\"{e}ren: de syntax is
%% \renewcommand{\filenaamlogo}{inputs/....}
%% waarbij je .... vervangt door de filenaam van het logo. Alvast drie beschikbare mogelijkheden zijn:
%% \renewcommand{\filenaamlogo}{inputs/kubus-wetenschappen.png}
%% \renewcommand{\filenaamlogo}{inputs/esat-logo.pdf}


\renewcommand{\filenaamlogo}{inputs/logo-innovationlab.png}

%% In de volgende lijn importeer je (in pdf) de frontpagina. Die maken we apart aan in een uniforme stijl. Er worden andere kleuren gebruikt voor de leerlingenbundels en voor de lerarenbundels.
%% Je kan de volgende regel veranderen in
%% \includepdf{inputs/voorbeeld-cover-vooraan-lkr.pdf}

\includepdf{inputs/frontcover.pdf}

%% Verander de volgende lijn niet.

\tableofcontents
%%%%%%%%%%%%%%%%%%%%%%%%%%%%%%%%%%%%%%%%%%%%%%%%%%%%%%%%%%%%%%%%%%%%%%%%%
%%%%%%%%%%%%%%%%%%%%%%%%%%%%%%%%%%%%%%%%%%%%%%%%%%%%%%%%%%%%%%%%%%%%%%%
%%%%%%%%%%%%%%%%%%%%%%%%%%%%%%%%%%%%%%%%%%%%%%%%%%%%%%%%%%%%%%%%%%%%%%%
\deelmetoef{Module 1}{Inleiding}{Module 1. Inleiding}{Oplossingen module 1}{Oplossingen module 1}
%%%%%%%%%%%%%%%%%%%%%%%%%%%%%%%%%%%%%%%%%%%%%%%%%%%%%%%%%%%%%%%%%%%%%%%
%%%%%%%%%%%%%%%%%%%%%%%%%%%%%%%%%%%%%%%%%%%%%%%%%%%%%%%%%%%%%%%%%%%%%%%

\begin{samenvatting}
Tijdens dit project maken we een smartphone app waarmee je je hartslag kan meten. Waarom het belangrijk is af en toe je hartslag te meten en waarom we dit met een app willen doen, wordt hierna uitgelegd. Daarna bekijken we wat een hartslag nu precies is, en hoe we die met onze smartphone kunnen meten.
\end{samenvatting}
%

%%%%%%%%%%%%%%%%%%%%%%%%%%%%%%%%%%%%%%%%%%%%%%%%%%%%%%%%%%%%%%%%%%%%%%%
\section{Wat en waarom?}
\label{sec:Mod1_Sec1}
%%%%%%%%%%%%%%%%%%%%%%%%%%%%%%%%%%%%%%%%%%%%%%%%%%%%%%%%%%%%%%%%%%%%%%%
%
De hartslag duidt aan hoe vaak je hart klopt per minuut. Tijdens het sporten is het belangrijk je hartslag in het oog te houden. Een te hoge hartslag wijst op een te hoge inspanning, en kan mogelijks gevaarlijk zijn. Bovendien kan een (te) hoge hartslag in rust wijzen op aandoeningen zoals een overactieve schildklier, bloedarmoede of bloedverlies, een longaandoening enz. Regelmatig je hartslag controleren is dus de boodschap.

Om je hartslag te meten wordt meestal gebruik gemaakt van een bloeddrukmeter.

\begin{minipage}{.5\linewidth}
	\gewonefiguur{width=\linewidth}{module1/patient}
\end{minipage} 
\begin{minipage}{.5\linewidth}
	\gewonefiguur{width=\linewidth}{module1/hartslagmonitor}
\end{minipage} 

Een bloeddrukmeter is een gespecialiseerd toestel dat meestal enkel een dokter of iemand met hartklachten ter beschikking heeft. Bovendien is het nogal omslachtig een bloeddrukmeter tijdens het sporten mee te nemen, gezien dit niet echt een draagbaar toestel is.

\gewonefiguur{width=.5\linewidth}{module1/lopermethartslagmeter}

Naast een bloeddrukmeter kan ook een hartslagmeter gebruikt worden om de hartslag te meten. 

\begin{minipage}{.5\linewidth}
	\gewonefiguur{width=\linewidth}{hartslagmeter}
\end{minipage} 
\begin{minipage}{.5\linewidth}
	\gewonefiguur{width=\linewidth}{module1/loper-met-dure-loopband}
\end{minipage} 

In tegenstelling tot een bloeddrukmeter kan je een hartslagmeter wel makkelijk meenemen bij het sporten en geeft ook real time feedback. Maar wellicht ligt de aanschaf van een hartslagmeter niet binnen het budget van elke leerling of student?

Daarom bouwen wij in dit project onze eigen smartphone app, waarmee we snel, eenvoudig en redelijk accuraat onze hartslag kunnen meten, zonder nood aan gespecialiseerde toestellen.

\gewonefiguur{width=.5\linewidth}{module1/loper}

\section{De hartslag}
\label{sec:Mod1_Sec2}

Je hart is een spier, die samentrekt om je bloed in je lichaam rond te pompen. De hartslag is de pompbeweging van je hart. In de volksmond is de hartslag ook het aantal samentrekkingen van je hart per minuut, of anders gezegd het aantal keer dat je hart klopt per minuut.

De hartslag van een volwassene varieert tussen 60 en 100 slagen per minuut. De maximale hartslag is de hoogste hartslag die je lichaam aankan gedurende fysische activiteit. De maximale hartslag varieert naargelang je leeftijd en kan berekend worden met de volgende vuistregel:

\begin{equation*}
MH = 220-LT
\end{equation*}

waarbij MH staat voor maximale hartslag (in slagen per minuut of \textquotedblleft beats per minute \textquotedblright (kortweg: \emph{bpm})) en LT voor je leeftijd (in \emph{jaar}).

Het wordt aangeraden tijdens het sporten je hartslag tussen 55 en 85\% van je maximale hartslag te houden. Voor een volwassene van 20 jaar oud is de maximale hartslag 200 bpm. Het aangeraden hartslagbereik tijdens het sporten is dus tussen $0,55*200$ en $0,85*200$ of nog tussen 110 en 170 bpm.

\begin{oef}
Bereken de maximale hartslag en het ideale hartslagbereik tijdens het sporten van een volwassene van 
\begin{itemize}
	\item 18 jaar
	\item 30 jaar
	\item 65 jaar
\end{itemize}
\end{oef}
\oplos{\begin{itemize}
		\item $MH_{18}=202$ en hartslagbereik tussen 111,1 en 171,7
		\item $MH_{30}=190$ en hartslagbereik tussen 104,5 en 161,5
		\item $MH_{65}=155$ en hartslagbereik tussen 85,25 en 131,75
	\end{itemize}
}

\begin{oef}Maak ook een grafiek die de maximale hartslag uitzet in functie van de leeftijd. 

Welke grootheid komt op de $x$-as? Wat is de eenheid? 

Welke grootheid komt op de $y$-as? Wat is de eenheid?
\end{oef}
\oplos{
	\begin{tikzpicture}[domain=0:100]
	\begin{axis}[axis lines=middle,xmax=110,ymin=-1,ymax=240,
		x label style={at={(axis description cs:0.5,-0.1)},anchor=north},
		y label style={at={(axis description cs:-0.1,.5)},rotate=90,anchor=south},
		xlabel={LT (Jaar)},ylabel={MH (bpm)}]
	\addplot[domain=0:100] {220-x} node[left]{$MH=220-LT$};
	\end{axis}
	\end{tikzpicture}
}

De hartslag, of het rondpompen van je bloed, voel je als je je hand op je hart legt. Maar ook in je hals of aan je pols kan je je hartslag voelen. De hartslag wordt dan ook vaak handmatig gemeten door twee gestrekte vingers op de pols (van jezelf of van een ander persoon) of in de hals (tegen de halsslagader) te plaatsen. Hier voel je de hartslag heel goed. Door gedurende 15 seconden het aantal hartslagen te tellen en te vermenigvuldigen met 4 kan je de hartslag, dit is het aantal keer dat je hart pompt per minuut, bepalen.

\begin{opdracht}{Je eigen hartslag meten}
	Meet je eigen hartslag door twee gestrekte vingers in je hals te plaatsen, het aantal keer dat je hart pompt gedurende 15 seconden te tellen en dit te getal te vermenigvuldigen met vier.
	
	Als je moeite hebt om het aantal hartslagen te tellen, kan je ook gebruik maken van de smartphone app \textquoteleft Science journal\textquoteright. Start een nieuw experiment en meet het omgevingsgeluid, zoals in de figuur hieronder. Plaats twee gestrekte vingers in je hals en zeg 'ja' elke keer als je een hartslag voelt. Doe dit gedurende 15 seconden. Tel nadien het aantal pieken die geregistreerd werden en vermenigvuldig dit aantal met vier om de hartslag te bekomen.\footnote{Tip voor leerkachten:
	Het gebruik van de app Science Journal kan je demonstreren door je smartphone scherm te projecteren op je computer. Hiervoor kan je het programma (op je pc) en de app (op je smartphone) Vysor gebruiken wanneer je smartphone met een USB kabel met de pc verbonden is.}
	
	%TODO hier grafiek van maken en invoegen als figuur
	
\end{opdracht}

\begin{oef}
	Met welke factor moet je het aantal hartslagen gemeten gedurende een telperiode van 10 seconden vermenigvuldigen om de hartslag in \emph{bpm} te bekomen? 
	
	En wat als de telperiode 20 seconden is?
	
	Wat is het voordeel van een langere telperiode?
\end{oef}
\oplos{\begin{itemize}
		\item $10~s$: factor $6$
		\item $20~s$: factor $3$
		\item Hoe korter de telperiode, hoe nauwkeuriger je moet tellen. Eventuele telfouten worden immers vermenigvuldigd met de factor. De factor is groter, naarmate de telperiode korter is. Telfouten worden dus groter als de telperiode korter is.
	\end{itemize}
}

Voor heel hoge hartslagen ($>160~bpm$) wordt het quasi onmogelijk de hartslag nog zelf te tellen. Dan moet je een meettoestel gebruiken. In wat volgt maken wij ons eigen meettoestel: een smartphone app waarmee je goedkoop, snel en gemakkelijk je hartslag kan meten.

\section{Overzicht van het project}
\label{sec:Mod1_Sec3}

Maar hoe kunnen we met onze smartphone onze hartslag gaan meten? We hebben al gezien dat je het pompen van je hart goed kan voelen aan je pols of in de hals. Net zo is de hartslag ook waarneembaar in je vingertop; ook daar zijn kleine pulsaties in je aders door het pompen van je hart. 

\gewonefiguur{height=10cm}{module1/loperStilstaand}

We zullen gebruik maken van de camera die in je smartphone ingebouwd is. We zullen onze vingertop zachtjes op de lens van de camera plaatsen. De kleine pulsaties in de aders van je vinger resulteren in kleine veranderingen in de kleur van je vinger. Met het blote oog kan je die miniscule kleurveranderingen niet waarnemen, maar gelukkig kan je die subtiele veranderingen wel met je smartphone detecteren. De kleurveranderingen in functie van de tijd zijn gerelateerd aan het pompen van je hart. Een grafiek van de kleurverandering in functie van de tijd zie je hieronder. Door opnieuw het aantal pieken in de telperiode van 15 seconden te tellen en te vermenigvuldigen met vier bekom je opnieuw de hartslag.

%TODO figuur maken en invoegen (Excel?)

De opeenvolging van uit te voeren stappen kunnen we voorstellen in een blokdiagram, zoals hieronder. Dit blokdiagram toont welke stappen we moeten doorlopen om de hartslag te bepalen, zonder in detail uit te werken wat precies moet gebeuren in elk blok.

%TODO blok schema verder uitleggen

\begin{tikzpicture}[thick]
\node[draw,rectangle] at (0,0) (a) {Video opnemen};
\node[draw,rectangle] at (0,-2) (b) {Video naar afbeeldingen};
\node[draw,rectangle] at (0,-4) (c) {Kleurwaarde};
\node[draw,rectangle] at (5,-1) (d) {Grafiek kleurwaarde};
\node[draw,rectangle,dashed] at (10,-1) (j) {Hartslag};
\node[draw,rectangle] at (5,-4) (e) {Piekdetectie};
\node[draw,rectangle] at (10,-4) (f) {Hartslag};
\node[draw,rectangle] at (5,-7) (g) {Fouriertransformatie (DFT)};
\node[draw,rectangle] at (10,-7) (h) {Grafiek DFT};
\node[draw,rectangle,dashed] at (13,-7) (i) {Hartslag};

\draw[->] (a) -- (b) ;
\draw[->] (b) -- (c) ;

\draw[->] (c) -- (d) ;
\draw[->] (c) -- (e) ;
\draw[->] (c) -- (g) ;

\draw[->,dashed] (d) -- (j) ;
\draw[->] (e) -- (f) ;
\draw[->] (g) -- (h) ;
\draw[->,dashed] (h) -- (i) ;

\end{tikzpicture}

%\begin{tikzpicture}
%\sbEntree{n1}
%\sbRelier[RecordButton]{n1}
%\sbBlocL{n1}{Video opnemen}{n2}
%%\sbBlocL{n2}{Video naar afbeeldingen}{n3}
%%\sbBlocL{n3}{Kleurwaarde van afbeeldingen}{n4}
%%\sbBlocL{n4}{Grafiek van kleurwaarde ifv tijd}{n5}
%%\sbBlocL{n5}{Piekdetectie kleurwaarden}{}
%%\sbBlocL{n6}{Hartslag berekenen}{n7}
%%\sbSortie[5]{S1}{n7}
%%\sbNomLien[0.8]{S1}{$Hartslag$}
%
%%\sbEntree{E}
%%\sbComp{a}{E}
%%\sbBloc{b}{$H_1$}{a}
%%\sbRelier[$E_1$]{E}{a}
%%\sbBlocL{c}{$H_2$}{b}
%%\sbRelier[$\epsilon$]{a}{b}
%%\sbComph{d}{c}
%%\sbRelier[u]{c}{d}
%%\sbBlocL{e}{$H_3$}{d}
%%\sbBlocL{f}{$H_4$}{e}
%%\sbSortie[5]{S1}{f}
%%\sbRelier{f}{S1}
%%\sbNomLien[0.8]{S1}{$S_1$}
%%\sbDecaleNoeudy[-4]{f}{u}
%%\sbDecaleNoeudy{e}{v}
%%\sbBlocr{r1}{$R_1$}{u}
%%\sbBlocr{r2}{$R_2$}{v}
%%\sbBlocrL{r3}{$R_3$}{r2}
%%\sbRelieryx{f-S1}{r1}
%%\sbRelierxy[n1]{r1}{d}
%%\sbRelieryx{e-f}{r2}
%%\sbRelierxy[n2]{r3}{a}
%\end{tikzpicture}





%%%%%%%%%%%%%%%%%%%%%%%%%%%%%%%%%%%%%%%%%%%%%%%%%%%%%%%%%%%%%%%%%%%%%%%
%%%%%%%%%%%%%%%%%%%%%%%%%%%%%%%%%%%%%%%%%%%%%%%%%%%%%%%%%%%%%%%%%%%%%%%
\deelmetoef{Module 2}{Apps maken met MIT App Inventor 2}{Module 2. Apps maken met MIT App Inventor 2}{Oplossingen module 2}{Oplossingen module 2}
%%%%%%%%%%%%%%%%%%%%%%%%%%%%%%%%%%%%%%%%%%%%%%%%%%%%%%%%%%%%%%%%%%%%%%%
%%%%%%%%%%%%%%%%%%%%%%%%%%%%%%%%%%%%%%%%%%%%%%%%%%%%%%%%%%%%%%%%%%%%%%%

\begin{samenvatting}
De hartslagmonitor die we gaan maken is een smartphone app. MIT, een bekende universiteit in de Verenigde Staten, ontwikkelde een programmeeromgeving waarmee je zeer eenvoudig apps kan maken en testen. In dit hoofdstuk bespreken we de belangrijkste structuren voor het maken van een app. Ook in andere programmeertalen en -omgevingen zal je deze structuren kunnen gebruiken.
\end{samenvatting}
%

%%%%%%%%%%%%%%%%%%%%%%%%%%%%%%%%%%%%%%%%%%%%%%%%%%%%%%%%%%%%%%%%%%%%%%%
\section{MIT App Inventor 2}
\label{sec:Mod2_Sec1}
%%%%%%%%%%%%%%%%%%%%%%%%%%%%%%%%%%%%%%%%%%%%%%%%%%%%%%%%%%%%%%%%%%%%%%%
%
Met behulp van MIT App Inventor 2 kan je zelf apps maken. Alles wat je nodig hebt, is een computer met wifi-verbinding en een smartphone die met hetzelfde wifi-netwerk geconnecteerd is.

\subsection{Maken van de app}

Surf naar de website \url{http://appinventor.mit.edu/explore/} en klik rechtsboven op create apps! Je zal moeten inloggen met een google account. Als je nog geen google account hebt, kan je die aanmaken via \url{https://accounts.google.com/signup}. Na succesvol aanmelden kom je op de pagina met je projecten terecht. 

Cre\"eer een nieuw project, noem het bv \textquotedblleft HelloWorld\textquotedblright, en klik erop om het project te openen. 

Je komt op het startscherm terecht. Er zijn twee \textquoteleft views\textquoteright. De eerste view is de \emph{designer view}, weergegeven in Figuur \ref{fig:view_designer}.

\figuurmetlabel[\label{fig:view_designer}]{width=.8\linewidth}{module2/overviewAI}{Designer view in MIT App Inventor 2. Overschakelen tussen views doe je rechtsbovenaan (blauwe rechthoek). Tabs zijn aangeduid in rood.}

Hierin kan je de componenten toevoegen die je zal zien op het scherm tijdens het uitvoeren van je app. In de tab \emph{Palette} (links) vind je alle componenten die je kan toevoegen. Dit is bijvoorbeeld een knop (Engels: button), een label, een afbeelding (Engels: image) enz. Je voegt componenten toe door ze vast te nemen en ze naar het scherm van de afgebeelde smartphone te slepen. Componenten kan je groeperen m.b.v. \texttt{HorizontalArrangement}, \texttt{VerticalArrangement} of \texttt{TableArrangement}. Componenten kan je een naam geven onder de tab \emph{Components} door de component aan te klikken en op \emph{Rename} te drukken. In de tab \emph{Properties} (rechts) kan je de eigenschappen van de component veranderen: Welke tekst moet op de knop staan? Wat is de lettergrootte? Wil je de tekst in een andere kleur weergeven? ... Als je wilt weten wat een component doet, kan je het vraagteken naast de component in de tab \emph{Palette} aanklikken voor een korte uitleg.

De andere view is de \emph{blocks view}, weergegeven in Figuur \ref{fig:view_block}.

\figuurmetlabel[\label{fig:view_block}]{width=.8\linewidth}{module2/uitlegAI_blocks}{Block view in MIT App Inventor 2.}

Hierin voeg je de functionaliteit toe aan je app. Wat moet er gebeuren als je op de knop drukt? Misschien wil je de tekst in een label aanpassen, of wil je een andere afbeelding tonen. Er kan veel; wat je implementeert, hangt af van je verbeelding.

\begin{opdracht}{Experimenteren met MIT App Inventor 2}
Ga na wat je allemaal kan doen in de programmeeromgeving. Welke componenten bestaan er allemaal? Welke eigenschappen kan je aanpassen?

Probeer een eerste app te maken. De app bevat een knop met de tekst \textquotedblleft Zeg hallo\textquotedblright, en een label zonder tekst. Als je op de knop drukt, moet de tekst \textquotedblleft Hello world!\textquotedblright \ in het label verschijnen. Dit \textquotedblleft Hello world!\textquotedblright-programma is een van de simpelste programmatjes die je kan schrijven. Wanneer programmeurs een nieuwe programmeertaal aanleren, schrijven ze steeds \textquotedblleft Hello world!\textquotedblright-programma om te leren werken met de nieuwe taal, die we met een geleerder woord syntax noemen.
\end{opdracht}

\subsection{Testen en uitvoeren van de app}
Om de app te testen moet je de MIT App Inventor 2 Companion App (te vinden in de Google Play Store) installeren op je smartphone. 
Zorg ervoor dat je computer en smartphone met hetzelfde wifi-netwerk verbonden zijn. 
Klik op connect, rechtsboven in Figuur \ref{fig:view_designer}, en kies de eerste optie \textquotedblleft AI companion\textquotedblright. Er verschijnt een QR-code die je kan scannen met je smartphone met de MIT App Inventor 2 Companion App. De app opent op je smartphone, zodat je de werking van de app kan testen.

Als je smartphone geconnecteerd is via de Companion App, verandert de app op je smartphone mee wanneer je aanpassingen doet aan de componenten op het scherm. Soms werkt dit echter niet zo goed. Breek dan de connectie met de smartphone af via \texttt{Connection - reset connection} en maak de verbinding opnieuw via \texttt{Connection - AI companion}.

Als je de app die je ontworp getest hebt via de App Inventor 2 Companion App en de app aan de vereisten voldoet, kan je de app downloaden. Dan hoef je niet meer steeds de QR code te scannen als je de app wilt gebruiken. De app kan je installeren via \emph{Build - App (provide QR code for .apk)}. Door een QR code te scannen met je GSM kan je het installatie-bestand (.apk) downloaden op je smartphone. Wanneer je het .apk-bestand opent, installeert de app.

\begin{opdracht}{Good morning app}
	Toon \textquotedblleft Good morning!\textquotedblright \ in drie verschillende talen op het scherm, nl. Italiaans, Spaans en Duits.
	
	\begin{minipage}{.5\linewidth}
		\gewonefiguur{width=5cm}{module2/goedemorgenApp}
	\end{minipage}
	\begin{minipage}{.5\linewidth}
		\begin{enumerate}
			\item Cre\"eer een nieuw project.
			\item Ontwikkel de designer view:
			\begin{itemize}
				\item Welke layout?
				\item Welke componenten?
				\item Welke naam voor elke component?
				\item Welke eigenschappen voor elke component?
			\end{itemize}
			\item Ontwikkel de block view
			\begin{itemize}
				\item Indien klik op knopItaliaans: \textquotedblleft Buongiorno!\textquotedblright
				\item Indien klik op knopSpaans: \textquotedblleft Buenos dias!\textquotedblright
				\item Indien klik op knopDuits: \textquotedblleft Guten Morgen!\textquotedblright
			\end{itemize}
		\end{enumerate}
	\end{minipage}

Test de werking van je app grondig. Doet de app op elk moment wat je ervan verwacht? Indien niet: ga op zoek naar de onderliggende oorzaak. Dit proces, waarin je ongewenste fouten in de werking van de app opspoort en wegwerkt, heet \textquotedblleft debuggen \textquotedblright. Zelfs de beste programmeurs moeten hun apps debuggen. Dit is soms een frustrerend proces, maar het is helaas wel een wezenlijk onderdeel van programmeren.

\end{opdracht}

%%%%%%%%%%%%%%%%%%%%%%%%%%%%%%%%%%%%%%%%%%%%%%%%%%%%%%%%%%%%%%%%%%%%%%%
\section{Blokken in de blocks view}
\label{sec:Mod2_Sec2}
%%%%%%%%%%%%%%%%%%%%%%%%%%%%%%%%%%%%%%%%%%%%%%%%%%%%%%%%%%%%%%%%%%%%%%%
%
Met MIT App Inventor 2 kan je heel gemakkelijk zelf aan de slag. De betekenis van de componenten is vaak intu\"itief duidelijk en via het vraagteken kan je extra uitleg oproepen. 

Hieronder bespreken we enkele belangrijke blokken uit de blocks view, waarvan de bedoeling misschien iets minder vanzelfsprekend is.

\subsection{Controle (control)}
De meeste programmeertalen zijn opgebouwd rond drie belangrijke controlestructuren: \textquoteleft if\textquoteright-loops om beslissingen te nemen, en \textquoteleft for\textquoteright-loops en \textquoteleft while\textquoteright-loops om herhaling in te bouwen.

\subsubsection{Beslissingen nemen: \emph{If then} blok}

\begin{minipage}{.5\linewidth}
\gewonefiguur{width=7cm}{module2/if.jpg}
\end{minipage}
\begin{minipage}{.5\linewidth}
\gewonefiguur{width=7cm}{module2/ifStructures}
\end{minipage}

\textquoteleft If\textquoteright-lussen laten toe een programma dynamisch te maken. Een \textquoteleft If\textquoteright-lus kan je vergelijken met het kiezen van een weg op een kruispunt: als een voorwaarde voldaan is (YES), ga je naar links; anders ga je rechts. 

Het testen of een voorwaarde voldaan is, gebeurt met een \texttt{Boolean} expressie.
Een \texttt{Boolean} expressie is een uitdrukking die ofwel waar of ofwel vals is.
Het resultaat van \texttt{Boolean} expressie is een \texttt{Boolean}, een variabele die maar twee waarden kan aannemen: \texttt{true} als de voorwaarde voldaan is, en \texttt{false} als de voorwaarde niet voldaan is.

Je kan kiezen om een bewerking enkel uit te voeren als een voorwaarde voldaan is (\emph{If then} blok) of de ene bewerking uit te voeren als de voorwaarde voldaan is, en een andere als de voorwaarde niet voldaan is (\emph{If then else} blok). Als je echt complexe dingen wil doen, kan je extra voorwaarden toevoegen (\emph{If then, else if then, else then} blok).

\subsubsection{Herhaling: \emph{While test} lus of \emph{for each} lus}

Voor zaken die vaak uitgevoerd moeten worden, bv. de tekst op een label laten optellen van 1 tot 100, kan je gebruik maken van een repetitie-blok, zoals een \emph{While test} lus of \emph{for each} lus.

\gewonefiguur{width=\linewidth}{module2/repititionStructures2}

Bij een \emph{for each}-lus wordt een teller (in ons voorbeeld \texttt{number}) steeds verhoogd (in ons voorbeeld met $Z$) binnen een bereik van waarden (in ons voorbeeld tussen $X$ en $Y$). Een \emph{for}-lus is ideaal om een probleem op te lossen waarbij een lus een vast aantal keer herhaald moet worden.

Een \emph{While test} lus heeft twee delen: (1) een \emph{test}-socket: \texttt{Boolean} expressie (waar of niet waar?) die getest wordt, en (2) een \emph{do}-socket: een verzameling instructies die uitgevoerd wordt zolang de \texttt{Boolean} expressie \texttt{true} is. Pas op dat je geen oneindige lussen maakt: ooit moet de \texttt{Boolean} expressie \texttt{false} worden, anders blijf je vastzitten in de \emph{while}-lus.

\subsection{Logica (Logic)}

Zoals eerder gezegd is een \texttt{Boolean} expressie een expressie die ofwel waar of ofwel vals is.
Logische operatoren (zoals \texttt{AND} en \texttt{OR}) kunnen gebruikt worden om meerdere boolean expressies samen te stellen en zo meer complexe voorwaarden te testen. 

De logische operator \texttt{AND} geeft als output waar als beide inputs waar zijn, en onwaar als een of meerdere inputs onwaar zijn. 

De logische operator \texttt{OR} geeft als output waar als minstens een van beide inputs waar is, en onwaar als beide inputs onwaar zijn.

De logische operator \texttt{NOT} geeft als output onwaar als de input waar is en geeft als output waar als de input onwaar is.

\subsection{Wiskunde (Math)}
Uiteraard kunnen bewerkingen op getallen uitgevoerd worden. Je vindt het volledige overzicht van mogelijke bewerkingen onder \emph{Blocks - Math}.

\subsection{Tekst (Text)}
Naast wiskundige bewerkingen kan je ook tekst input van de gebruiker verwerken. De mogelijke bewerkingen die je op tekst kan uitvoeren vind je onder \emph{Blocks - Text}.

\subsection{Variabelen (Variables)}
Variabelen worden gebruikt om waarden bij te houden. Een bij te houden waarde kan bv. het resultaat van een wiskunde bewerking zijn, of een tekst input van de gebruiker die je later nog wilt aanpassen. 

Belangrijk is dat er een onderscheid gemaakt wordt tussen lokale en globale variabelen, zie Figuur \ref{fig:variabele}. Globale variabelen worden aangemaakt bij de opstart van het programma en kunnen doorheen het hele programma gebruikt worden. We zouden kunnen zeggen dat je vanuit elk blok \textquotedblleft aan de globale variabele kan\textquotedblright.
Lokale variabelen worden aangemaakt binnen een blok en kunnen enkel door blokken binnen de grenzen van de lokale variabele gebruikt worden. Je \textquotedblleft kan er dus enkel aan\textquotedblright \ binnen het oranje blok zelf. Vandaar dat het blok voor een lokale variabele onderaan nog een oranje staartje heeft. Dit staartje duidt aan binnen welke grenzen \textquotedblleft je aan de variabele kan\textquotedblright.

\figuurmetlabel[\label{fig:variabele}]{width=8cm}{module2/variabelen}{Het bovenste blok toont hoe je een globale variabele kan aanmaken, die doorheen het hele programma gebruikt kan worden. Het onderste blok toont hoe je een lokale variabele kan aanmaken, die enkel toegankelijk is binnen het blok van de variabele.}

De \texttt{Get}-instructies worden gebruikt om de waarde van een variabele op te halen. De \texttt{Set}-instructie wordt gebruikt om de waarde van een variabele te wijzigen.

\subsection{Lijsten (Lists)}
Een lijst kan meerdere items met gelijkaardige data bevatten. Zo kan je bijvoorbeeld een lijst van getallen hebben, of een lijst van tekstwaarden. Als je 100 getallen moet bijhouden, is het makkelijker een lijst te gebruiken dan 100 verschillende variabelen.

Nadat een lijst wordt aangemaakt zoals in Figuur \ref{fig:lijst}, kan je extra items aan de lijst toevoegen, of er items uit verwijderen. De volledige lijst van bewerkingen die je op een lijst kan toepassen, vind je onder \emph{Blocks - list}

\figuurmetlabel[\label{fig:lijst}]{width=8cm}{module2/lijst}{Je kan een lege lijst aanmaken (boven), of een lijst aanmaken die reeds een aantal elementen bevat.}



%%%%%%%%%%%%%%%%%%%%%%%%%%%%%%%%%%%%%%%%%%%%%%%%%%%%%%%%%%%%%%%%%%%%%%%
%%%%%%%%%%%%%%%%%%%%%%%%%%%%%%%%%%%%%%%%%%%%%%%%%%%%%%%%%%%%%%%%%%%%%%%
\deelmetoef{Module 3}{Digitale afbeeldingen en video}{Module 3. Digitale afbeeldingen en video}{Oplossingen module 3}{Oplossingen module 3}
%%%%%%%%%%%%%%%%%%%%%%%%%%%%%%%%%%%%%%%%%%%%%%%%%%%%%%%%%%%%%%%%%%%%%%%
%%%%%%%%%%%%%%%%%%%%%%%%%%%%%%%%%%%%%%%%%%%%%%%%%%%%%%%%%%%%%%%%%%%%%%%

\begin{samenvatting}
Het meten van je hartslag met de smartphone is gebaseerd op kleurveranderingen in je vingertop, die we met de smartphone camera registreren in een video. Om goed te begrijpen hoe we die kleurveranderingen precies kunnen detecteren, is het belangrijk de basisprincipes achter digitale afbeeldingen en video te begrijpen. Daar gaan we in deze module dieper op in.
\end{samenvatting}
%

%%%%%%%%%%%%%%%%%%%%%%%%%%%%%%%%%%%%%%%%%%%%%%%%%%%%%%%%%%%%%%%%%%%%%%%
\section{Digitale afbeeldingen}
\label{sec:Mod3_Sec1}
%%%%%%%%%%%%%%%%%%%%%%%%%%%%%%%%%%%%%%%%%%%%%%%%%%%%%%%%%%%%%%%%%%%%%%%
%

\gewonefiguur{width=4cm}{module3/qrcode_youtubespreadsheets}

\begin{itemize}
	\item Matrices RGB - afbeelding Dimitri
	\item Filmpje spreadsheets \url{https://www.youtube.com/watch?v=UBX2QQHlQ_I&t=213s}
	\item Roodwaarde berekenen in App Inventor
\end{itemize}

%%%%%%%%%%%%%%%%%%%%%%%%%%%%%%%%%%%%%%%%%%%%%%%%%%%%%%%%%%%%%%%%%%%%%%%
\section{Digitale video}
\label{sec:Mod3_Sec2}
%%%%%%%%%%%%%%%%%%%%%%%%%%%%%%%%%%%%%%%%%%%%%%%%%%%%%%%%%%%%%%%%%%%%%%%
%

\gewonefiguur{width=4cm}{module3/qrcode_youtubeDigitalVideo}

\begin{itemize}
	\item Uitleg video - afbeelding Dimitri
	\item Filmpje digital video \url{https://www.youtube.com/watch?v=-1s-SuUQYs4}
	\item App video to frames
\end{itemize}

%%%%%%%%%%%%%%%%%%%%%%%%%%%%%%%%%%%%%%%%%%%%%%%%%%%%%%%%%%%%%%%%%%%%%%%
\section{Toepassing: de hartslagmonitor}
\label{sec:Mod3_Sec3}
%%%%%%%%%%%%%%%%%%%%%%%%%%%%%%%%%%%%%%%%%%%%%%%%%%%%%%%%%%%%%%%%%%%%%%%
%

\begin{itemize}
	\item Gemiddelde roodwaarde voor filmpje naar excel bestand
	\item Intermezzo: App inventor 2: for, file om bestand op te slaan
	\item Hartslag bepalen in excel
\end{itemize}


%%%%%%%%%%%%%%%%%%%%%%%%%%%%%%%%%%%%%%%%%%%%%%%%%%%%%%%%%%%%%%%%%%%%%%%
%%%%%%%%%%%%%%%%%%%%%%%%%%%%%%%%%%%%%%%%%%%%%%%%%%%%%%%%%%%%%%%%%%%%%%%
\deelmetoef{Module 4}{Automatische piekdetectie}{Module 4. Automatische piekdetectie}{Oplossingen module 4}{Oplossingen module 4}
\label{sec:piekdetectie}
%%%%%%%%%%%%%%%%%%%%%%%%%%%%%%%%%%%%%%%%%%%%%%%%%%%%%%%%%%%%%%%%%%%%%%%
%%%%%%%%%%%%%%%%%%%%%%%%%%%%%%%%%%%%%%%%%%%%%%%%%%%%%%%%%%%%%%%%%%%%%%%

\begin{samenvatting}
	In de vorige modules leerde je werken met App Inventor 2 en leerde je de basis van digitale afbeeldingen en video. In de slotopdracht van module 3 maakte je een grafiek waaruit je je hartslag kon afleiden door zelf pieken te tellen. In deze module gaan we na hoe de smartphone zelf die pieken kan detecteren, zodat wij van dat rekenwerk gespaard blijven. Daarvoor staan we stil bij wat een piek precies is, en hoe we de smartphone kunnen leren pieken te detecteren. 
\end{samenvatting}
%

%%%%%%%%%%%%%%%%%%%%%%%%%%%%%%%%%%%%%%%%%%%%%%%%%%%%%%%%%%%%%%%%%%%%%%%
\section{Wat is een piek in een signaal?}
\label{sec:Mod4_Sec1}
%%%%%%%%%%%%%%%%%%%%%%%%%%%%%%%%%%%%%%%%%%%%%%%%%%%%%%%%%%%%%%%%%%%%%%%
%
I.p.v. zelf pieken te tellen op een grafiek willen we dit automatisch doen in de app. Dit houdt in dat we de code moeten schrijven om pieken te detecteren. Door onze ervaringen in het dagelijkse leven kunnen wij mensen redelijk intu\"itief pieken herkennen. De smartphone daarentegen moet door ons gezegd worden wanneer hij een extra piek moet tellen. Daarom moeten we zelf ook even stilstaan bij wat een piek nu eigenlijk is.

\begin{steroef}
	Omschrijf in eigen woorden wat een piek is.
\end{steroef}
\oplos{Wellicht komen hier heel diverse antwoorden, afhankelijk van de achtergrond van de leerlingen. Leerlingen uit de 3e graad zullen wellicht over maxima spreken, terwijl leerlingen uit de 2e graad misschien meer woordelijke/minder wiskundige omschrijvingen zullen geven.}

Een woordelijke omschrijving van een piek hebben we gevonden, maar je smartphone redeneert wiskundiger. 

\begin{oef}
	Stel een wiskundige uitdrukking op die je in de code kan gebruiken om pieken te detecteren. 
	
	Hou er rekening mee dat je smartphone enkel een lijst met roodwaarden als input heeft. Die lijst kunnen we voorstellen als een verzameling $\{ u_1, u_2, u_3, \ldots, u_n \}$ met $n$ elementen, waarbij $u_i$ de $i$-de berekende gemiddelde roodwaarde voorstelt.
	
	Als je dit niet meteen ziet, kan je zelf een aantal voorbeeldlijsten maken en a.d.h.v. de vooorbeelden nagaan wanneer het signaal een piek bereikt.
\end{oef}
\oplos{Er is een maximum voor element $u_i$ als en slechts als: $u_i > u_{i-1}$ en $u_{i} < u_{i+1}$.}

\begin{opdracht}{Automatische piekdetectie in App Inventor 2}
Schrijf een app die automatisch pieken detecteert in een lijst. Hou bij (1) op welk volgnummer $i$ een piek $u_i$ gedetecteerd wordt, (2) hoeveel pieken in totaal gedetecteerd worden. Gebruik labels om die resultaten op je scherm weer te geven ter controle.

\begin{enumerate}
	\item Gebruik eerst een lijst die je zelf ingeeft met \emph{Lists - make a list}. Worden de pieken gedetecteerd op de plek die je verwacht? Pas je code aan tot je het verwachte resultaat bekomt.
	
	\item Gebruik nu een lijst met berekende roodwaarden, die we in de vorige module hebben opgeslaan in een CSV-bestand. Voeg daarom een \emph{File}-component toe je aan je app. Wanneer een knop \emph{Read CSV} ingedrukt wordt, moet het CSV-bestand die je opgeslaan hebt in de vorige module ingelezen worden. Kijk indien nodig even terug in je vorige project waar en onder welke naam je het CSV-bestand hebt opgeslaan. Als het CSV-bestand is ingelezen, moet de inhoud omgezet worden in een lijst en in een lijst geplaatst worden. Geef die lijst weer in een label en vergelijk de tekst met het CSV-bestand op een computer om te controleren dat dit correct gebeurde. 
	
	\item Gebruik de code uit stap 1 om automatisch pieken te detecteren in de lijst met roodwaarden. Komt dit overeen met het aantal pieken dat jij telt. Waarom (niet)?
	
	\begin{opmerking}
		Afhankelijk van het CSV-bestand dat je inleest, kan het zijn dat de lijst met roodwaarden in plaats van getallen als elementen lijsten met 1 element als elementen bevat. Om aan de roodwaarde te raken, moet je dan het eerste element halen uit elke lijst (die ook weer een element van de oorspronkelijke lijst is).
		
	%TODO indien onduidelijk tekening van Dimitri invoegen
	\end{opmerking}
	
	\begin{opmerking}
		Vaak is de lijst met roodwaarden een beetje \textquotedblleft ruizig \textquotedblright. Dit zorgt ervoor dat de app meer pieken detecteert dan wij zouden doen. Soms zien we een klein piekje in het signaal, dat we negeren, omdat het maar een kleine piek is, of omdat de piek te snel na een grote hartslagpiek komt. 
	\end{opmerking}

	\item We moeten de app dus leren dat sommige pieken moeten genegeerd worden, wanneer de piek niet echt hoog is, of heel snel komt. Hoe zou jij dit implementeren? Je hoeft dit niet meteen te doen, maar zorg wel dat je een plan hebt.
	
	\item De makkelijkste manier om de kleine piekjes te negeren is door de lijst met roodwaarden een beetje op te kuisen. De lijst met roodwaarden noemen we het signaal. De ruis die aanwezig is, is ongewenst. De ruis verandert vaak sneller dan het signaal zelf. Daardoor kunnen we de ruis verminderen door die snelle veranderingen weg te werken. Een eenvoudige manier om snelle veranderingen weg te werken is door opeenvolgende roodwaarden uit te middelen. Dit wordt ook wel laagdoorlaatfiltering genoemd: de snelle veranderingen, vaak ook hoge frequenties genoemd, worden weggefilterd uit het signaal.
	
	Schrijf code om opeenvolgende roodwaarden uit te middelen. Maak een nieuwe lijst, bv \emph{smoothRedValuesList}, waarin elk element een gemiddelde is van vier opeenvolgende elementen uit de originele lijst. Herinner je dat we een gemiddelde berekenen als een som van elementen, gedeeld door het aantal elementen dat opgeteld wordt.
	Voor het gemiddelde van 4 willekeurige getallen:
	\begin{equation*}
	\overline{a_4} = \text{Gemiddelde van $a_1$, $a_2$, $a_3$, $a_4$} = \frac{a_1+a_2+a_3+a_4}{n} = \frac{\sum_{i=1}^{4} a_i}{4}
	\end{equation*}
	
	\begin{opmerking}
		Tip: gebruik een \emph{for}-loop om de elementen van de nieuwe lijst \emph{smoothRedValuesList} te berekenen.
	\end{opmerking}

	\item Gebruik de code uit stap 1 om de pieken in de nieuwe lijst \emph{smoothRedValuesList} te berekenen. Bekom je een aanvaardbaar resultaat? Is het nodig meer of minder elementen van de originele lijst uit te middelen om tot een aanvaardbaar?
	
	\begin{opmerking}
	We spreken hier over een aanvaardbaar resultaat en niet over een correct resultaat. Het is goed mogelijk dat het aantal pieken dat je code detecteert licht afwijkt van het aantal pieken dat jij telt. Hoe sterk de afwijking is, hangt af van de kwaliteit van de opgenomen video. In Module \ref{sec:digafbvid} gingen jullie al na welke factoren invloed hebben op de kwaliteit van de video.
	\end{opmerking}

	\item Als je tevreden bent over de automatische piekdetectie, kan je het aantal pieken gedetecteerd tijdens de duur van de video herleiden naar een waarde voor de hartslag, het aantal pieken per minuut.
	
	Gebruik de frame rate en het aantal elementen in de lijst met roodwaarden om de duur van de video te berekenen. Gebruik de regel van drie om het aantal pieken gedurende de videoduur te herleiden naar het aantal pieken per seconde. Dit is de gemeten hartslag. Geef die waarde weer in een label op het scherm. Bekom je een realistische waarde voor de gemeten hartslag?
\end{enumerate}
\end{opdracht}

%%%%%%%%%%%%%%%%%%%%%%%%%%%%%%%%%%%%%%%%%%%%%%%%%%%%%%%%%%%%%%%%%%%%%%%
\section{Intermezzo: differentiequoti\"ent}
\label{sec:Mod4_Sec2}
%%%%%%%%%%%%%%%%%%%%%%%%%%%%%%%%%%%%%%%%%%%%%%%%%%%%%%%%%%%%%%%%%%%%%%%
%
Bedoeld voor leerlingen uit de 3e graad.

Voor de bepaling of een piek in het signaal bereikt wordt, maakte je gebruik van een wiskundige uitdrukking gerelateerd aan het differentiequoti\"ent.

Stel dat $\{ u_1, u_2, \ldots, u_n\}$ een signaal is, bv. een lijst met roodwaarden. 

\begin{opmerking}
	We noemen dit een \emph{discreet} signaal (in de figuur hieronder in rood aangeduid): de meetwaarden (bv roodwaarden) zijn slechts op bepaalde tijdstippen bepaald. 
	
	Naast \emph{discrete} signalen bestaan ook \emph{continue} signalen (in de figuur hieronder in zwart aangeduid), waarvoor op elk moment een meetwaarde beschikbaar is. Het signaal loopt dus - zoals de naam het zegt - continu door.
	
	\gewonefiguur{width=.5\linewidth}{module4/discreetvscontinu}
	
	Vaak spreekt men ook over \emph{digitale} en \emph{analoge} signalen. \emph{Digitale} signalen zijn vergelijkbaar met \emph{discrete} signalen, terwijl \emph{analoge} signalen te vergelijken zijn met \emph{continue} signalen, hoewel niet exact dezelfde zaken bedoeld worden met de termen \emph{digitaal/discreet} en \emph{analoog/continu}.

\end{opmerking}

\begin{oef}
	In de lessen wiskunde leerde je al over functies. Een grafiek van de functie met functievoorschrift $y=f(x)$ toont hoe de grootheid $y$ verandert i.f.v. de grootheid $x$. Vaak kan de grootheid $x$ een continu bereik van waarden aannemen. In dat geval spreken we over een discreet/continu signaal. \emph{(Schrap wat niet past.)}
\end{oef}
\oplos{Als de grootheid $x$ een continu bereik van waarden kan aannemen, spreken we over een continu signaal.}

De signalen waar we hier mee werken, kan je dus enigszins vergelijken met functies die je kent uit de wiskundelessen. 
Het verschil is dat we hier slechts op discrete momenten een functiewaarde kunnen bepalen, terwijl in de wiskundeles meestal met continue signalen gewerkt wordt.

We gaan terug naar ons discreet signaal $\{ u_1, u_2, \ldots, u_n\}$. De uitdrukking die je kan om na te gaan of een piek bereikt was is:

\begin{equation*}
u_i-u_{i-1}
\end{equation*}

Immers, als $u_{i}-u_{i-1}>0$ en $u_{i+1}-u_{i}<0$ wordt een lokaal maximum, of een piek, bereikt op het tijdstip $i$.

Deze uitdrukking is sterk gerelateerd aan het differentiequoti\"ent $D$ voor een functie $f$:

\begin{equation}
D = \frac{f(a+\Delta h) - f(a)}{\Delta h}
\end{equation}

Voor zeer kleine waarden van $\Delta h$ benadert het differentiequoti\"ent de richtingsco\"effici\"ent van de raaklijn in het punt $(a,f(a))$. We zeggen ook weleens dat de afgeleide van de functie $f$ in het punt $(a,f(a))$ gelijk is aan 

\begin{eqnarray*}
f'(a) &=& \lim\limits_{\Delta h \to 0} D \\
&=& \lim\limits_{\Delta h \to 0} \frac{f(a+\Delta h)-f(a)}{\Delta h}
\end{eqnarray*}

Uit de wiskundelessen weet je dat een functie een extremum (= een maximum of een minimum) bereikt waar de afgeleide nul wordt. Om een \emph{minimum} te vinden zoeken we voor welke waarde van $x$ de functie overgaat van dalend (afgeleide $f'(x)< 0$) naar dalend (afgeleide $f'(x)> 0$). Om een \emph{minimum} te vinden zoeken we voor welke waarde van $x$ de functie overgaat van stijgend (afgeleide $f'(x)>0$) naar dalend (afgeleide $f'(x)<0$).

Wat we in dit project doen, is zeer gelijkaardig: we kijken voor welke waarde van $i$ $u_{i}-u_{i-1}$ overgaat van een positieve naar een negatieve waarde om een piek, dit is een maximum, te detecteren. 

\begin{oef}
	Nu we achterhaald hebben hoe onze aanpak in dit project vergelijkbaar is met afgeleiden en extrema-vraagstukken, kunnen we ons ook afvragen welke verschillen bestaan. Som op welke verschillen je ziet.
\end{oef}
\oplos{\begin{itemize}
		\item Bij afgeleiden gebruiken we continue signalen, en hier discrete.
		\item Bij afgeleiden is er nog een deling door $\Delta h$. Dit laten we hier achterwege. 
		 \begin{opmerking}
		 	Merk op dat je ook de uitdrukking $u_{i}-u_{i-1}$ zou kunnen delen door de tijdsduur tussen 2 foto's. Die tijdsduur is $\Delta i=1/fs$, het omgekeerde van de frame rate. 
		 	
		 	Die deling heeft geen invloed op het resultaat (de gedetecteerde pieken), omdat we zoeken voor welke waarde van $i$ de uitdrukking $u_{i}-u_{i-1}$ of $\frac{u_{i}-u_{i-1}}{\Delta i} = [ u_{i}-u_{i-1} ] fs$  overgaat van een positieve naar een negatieve waarde. Die extra (positieve) factor $fs$ be\"invloedt het teken, en dus ook de uitkomst, niet.
		 \end{opmerking}
	\end{itemize}}

%%%%%%%%%%%%%%%%%%%%%%%%%%%%%%%%%%%%%%%%%%%%%%%%%%%%%%%%%%%%%%%%%%%%%%%
\section{Intermezzo: Frequentieanalyse}
\label{sec:Mod4_Sec3}
%%%%%%%%%%%%%%%%%%%%%%%%%%%%%%%%%%%%%%%%%%%%%%%%%%%%%%%%%%%%%%%%%%%%%%%
%
Bedoeld voor leerlingen uit de 3e graad.

Je zal al ondervonden hebben dat de nauwkeurigheid van de hartslagmeting afhangt van de kwaliteit van het opgenomen filmpje. Als je vingertop of de smartphone bewoog tijdens de opname, is de grafiek van de roodwaarde in functie van de tijd ruiziger. Hoe ruiziger het signaal, hoe moeilijker de detectie van pieken wordt. En hoe moeilijker de piekdetectie, hoe meer kans op fouten en hoe onnauwkeuriger de meting.

Er bestaat een andere manier om de hartslag te bepalen, die niet gebaseerd in op piekdetectie en daardoor iets nauwkeuriger is. Deze alternatieve manier is gebaseerd op frequentieanalyse.

\subsection{Geluid: zuivere en samengestelde tonen}

We nemen even afstand van onze hartslagmonitor. Om goed te kunnen uitleggen wat we met frequentieanalyse bedoelen, kijken we even naar geluid.

Geluid is een kleine verandering in luchtdruk, die zich door de lucht voortplant. Elk luchtdeeltje ondervindt een trilling en geeft die trilling door aan de luchtdeeltjes in zijn buurt. De trilling van een luchtdeeltje in functie van de tijd kan er uitzien als volgt:

\gewonefiguur{width=\linewidth}{module4/sinusTrilling}

De trilling hierboven is een harmonische trilling, die als een sinusvormig signaal voorgesteld wordt en met als functievoorschrift 

\begin{equation*}
y(t) = Asin(2\pi f t)
\end{equation*}. 

Links is de trilling voorgesteld in het tijdsdomein: de uitwijking in functie van de tijd is een sinusoidaal verloop. Rechts is de trilling voorgesteld in het frequentiedomein. Deze voorstelling noemen we dan ook het \emph{frequentiespectrum}. De frequentie van een trilling is gedefineerd als het aantal trillingen per seconde. De eenheid is $1/s$ of Hertz ($Hz$). Hoe het frequentiespectrum gevonden wordt, bespreken we later.

Deze trilling is een \emph{zuivere toon}: er is slechts 1 frequentie aanwezig. In het frequentiespectrum zien we maar 1 lijn. Het is het geluid dat je hoort als je een stemvork aanslaat.

\begin{oef}
\begin{itemize}
	\item Wat is de uitwijking, of amplitude $A$ van de trilling? 
	\item Hoe verandert het geluid als de amplitude verhoogt?
	\item Wat is de frequentie $f$ van de trilling? Vergeet de eenheid niet!
	\item Hoe verandert het geluid als de frequentie verhoogt?
	\item De periode is $T=1/f$ is het tijdsverschil tussen 2 opeenvolgende maximale (of minimale) uitwijkingen. Het is dus de tijd die 1 volledige trilling inneemt.
\end{itemize}
\end{oef}

Naast zuivere tonen bestaan ook samengestelde tonen. De meeste geluiden die je kent, zijn samengestelde tonen, waarin meerdere frequenties voorkomen. Hieronder zie je een voorbeeld.

\gewonefiguur{width=\linewidth}{module4/samengesteldeToon}

Het functievoorschrift van de afgebeelde trilling is 

\begin{equation*}
y(t) = 0.5 sin(2 \pi 10 t) + 1 sin(2 \pi 20 t) + 1 sin(2 \pi 30 t).
\end{equation*}

De amplitude en frequentie van de samengestelde toon kan je uit het functievoorschrift afleiden: $A_1 = 0.5$, $f_1 = 10$ Hz, $A_2 = 1$, $f_2 = 20$ Hz, $A_3 = 1$ en $f_3 = 30$ Hz. Uit de linkergrafiek in het tijdsdomein kunnen die grootheden echter niet zo eenvoudig afgeleid worden. Gelukkig toont het frequentiespectrum ons wel de nodige informatie: er zijn drie zuivere trillingen, \'e\'en op 10 Hz met amplitude 0.5, \'e\'en op 20 Hz met amplitude 1 en \'e\'en op 30 Hz met amplitude 1.

In het algemeen kan een samengestelde toon geschreven worden als een som van $N$ zuivere tonen, elk met amplitude $A_n$ en frequentie $f_n$:

\begin{equation*}
y(t) = \sum_{n=1}^{N} A_n sin(2 \pi f_n t).
\end{equation*}

\subsection{Frequentieanalyse: het idee van Fourier}

Meer algemeen kan je om het even welk periodiek signaal opbouwen door sinussen met veelvouden van een bepaalde grondfrequentie $f_0$ en verschillende amplituden op te tellen. Naarmate je meer termen optelt, krijgt het signaal een strakkere vorm.

Je kan elk periodiek signaal dus schrijven als een reeks

\begin{equation*}
y(t) = a_0 + \sum_{k=1}^{\inf} (a_k cos(2 \pi k f_0 t) + b_k sin(2 \pi k f_0 t)).
\end{equation*}

Frequentieanalyse of fourieranalyse is het ontleden van een oneindig lang periodiek signaal in de verschillende frequentiecomponenten waaruit het bestaat. 

Fourieranalyse bestaat er dus in de co\"effici\"enten $a_0$, $a_k$ en $b_k$ te zoeken, die bij elke frequentie $k f_0$ horen.

%TODO uitleg hoe dit te berekenen. Halen uit JC CI.

Fourieranalyse op geluidsignalen toont ons welke frequenties in een geluid voorkomen.

\subsection{Toepassing: fourieranalyse voor onze hartslagmonitor}

Terug naar onze hartslagmonitor. Wij zijn op zoek naar onze hartslag: het aantal keer dat ons hart klopt per minuuut. Dit is eigenlijk ook een soort van frequentie: het kloppen van ons hart kunnen we vergelijken met een geluidstrilling. Elke keer ons hart samentrekt, bevindt het bloed in onze vingertop zich iets verder van de huid en is de roodwaarde in onze video minimaal. En elke keer ons hart maximaal uitgezet is, bevindt het bloed in onze vingertop zich iets dichter bij huid en is de roodwaarde in onze video maximaal. Onze hartslag is dus de frequentie waarmee ons hart klopt en is meteen ook de frequentie waarmee de roodwaarde in de opgenomen video fluctueert. Als we dus het frequentiespectrum van ons roodwaarde-signaal kunnen opstellen, kunnen we uit de pieken hierin de hartslag afleiden.

\begin{opdracht}{Fourieranalyse in App Inventor 2}
Gebruik Fourieranalyse om de hartslag in een roodwaarde signaal te bepalen.
	
\begin{itemize}
	
	\item Maak een nieuw project \texttt{IL\_hartslagmonitor\_fourieranalyse}.
	
	\item Voorzie een aantal knoppen: \emph{readCSV}, \emph{calculateFFT} en \emph{plotFFT}. Voorzie ook een textbox waarin de gebruiker de frame rate ($fs$) kan aanpassen en een knop \emph{fsChanged} waarmee de gebruiker kan aangeven dat hij de frame rate aangepast heeft. Als de knop \emph{fsChanged} aangeklikt wordt, moet een globale variabele aangepast worden aan de inhoud van de textbox.
	
	\item Gebruik nu een lijst met berekende roodwaarden, die we in de vorige module hebben opgeslaan in een CSV-bestand. Gebruik zoveel mogelijk code uit vorige opdrachten. Voeg daarom een \emph{File}-component toe je aan je app. Wanneer een knop \emph{Read CSV} ingedrukt wordt, moet het CSV-bestand die je opgeslaan hebt in de vorige module ingelezen worden. Kijk indien nodig even terug in je vorige project waar en onder welke naam je het CSV-bestand hebt opgeslaan. Als het CSV-bestand is ingelezen, moet de inhoud omgezet worden in een lijst en in een lijst geplaatst worden. Geef die lijst weer in een label en vergelijk de tekst met het CSV-bestand op een computer om te controleren dat dit correct gebeurde. 
	
	\item Kopieer de code uit een van de vorige opdrachten om opeenvolgende roodwaarden uit te middelen. Maak een nieuwe lijst, bv \emph{smoothRedValuesList}, waarin elk element een gemiddelde is van vier opeenvolgende elementen uit de originele lijst. 
	
	\begin{opmerking}
		Afhankelijk van het CSV-bestand dat je inleest, kan het zijn dat de lijst met roodwaarden in plaats van getallen als elementen lijsten met 1 element als elementen bevat. Om aan de roodwaarde te raken, moet je dan het eerste element halen uit elke lijst (die ook weer een element van de oorspronkelijke lijst is).
		
		%TODO indien onduidelijk tekening van Dimitri invoegen
	\end{opmerking}
	
 	\item De implementatie van de Fourieranalyse hebben we al voor jullie gedaan. Voeg daarom de extensie \emph{FFTCalculator1} toe aan je app.
 	
	\item Gebruik het blok \emph{FFTCalculator - CalculateFft} om een Fourieranalyse uit te voeren wanneer de knop \emph{calculateCSV} aangeklikt wordt. Het resultaat is een lijst \emph{redValuesFFT}. 
	
	\item Verwijder het eerste element uit de lijst. Dit is nodig om een correct frequentiespectrum te kunnen maken.
	
	\item Wanneer op de knop \emph{plotFFT} gedrukt wordt, moet het frequentiespectrum getoond worden. Gebruik hiervoor het blok \emph{Chartmaker - DrawLineGraph}. Voeg opnieuw een \emph{Webviewer} toe aan je app. De titel van de grafiek en van de assen en het label kies je zelf. Maak de grafiek eerst met \emph{scalingFactor} gelijk aan 1.
	
	\item Op de $x$-as staan de mogelijke hartslagfrequenties. De sprong tussen twee opeenvolgende hartslagfrequenties is $\frac{60 \cdot fs}{\text{aantal roodwaarden}}$. Pas de \emph{scalingFactor} hieraan aan.
	
	\item Welke hartslag meten we nu op? De grootste piek in de grafiek toont welke frequentie \textquotedblleft het meeste in het signaal aanwezig is \textquotedblright. De $x$-waarde van de grootste piek toont dus de meest waarschijnlijke hartslag. Je kan de $x$- en $y$- waarde van een punt op de grafiek zien door erop te klikken. Komt de hoogste piek overeen met de hartslag die je zelf berekende? En met de hartslag die je app berekende via automatische piekdetectie?
	
	\begin{opmerking}
		Het is mogelijk dat de grootste piek zichtbaar is voor een heel kleine $x$-waarde. In dat geval is de piek wellicht veroorzaakt door ruis in het signaal. Komt de 2e grootste piek wel in de buurt bij een meer waarschijnlijke hartslag?
	\end{opmerking}
\end{itemize}
\end{opdracht}







%%%%%%%%%%%%%%%%%%%%%%%%%%%%%%%%%%%%%%%%%%%%%%%%%%%%%%%%%%%%%%%%%%%%%%%
%%%%%%%%%%%%%%%%%%%%%%%%%%%%%%%%%%%%%%%%%%%%%%%%%%%%%%%%%%%%%%%%%%%%%%%
\deelmetoef{Module 5}{Bouwen van de complete app}{Module 5. Bouwen van de complete app}{Oplossingen module 5}{Oplossingen module 5}
%%%%%%%%%%%%%%%%%%%%%%%%%%%%%%%%%%%%%%%%%%%%%%%%%%%%%%%%%%%%%%%%%%%%%%%
%%%%%%%%%%%%%%%%%%%%%%%%%%%%%%%%%%%%%%%%%%%%%%%%%%%%%%%%%%%%%%%%%%%%%%%

\begin{samenvatting}
\end{samenvatting}
%

%%%%%%%%%%%%%%%%%%%%%%%%%%%%%%%%%%%%%%%%%%%%%%%%%%%%%%%%%%%%%%%%%%%%%%%
\section{Video opnemen}
\label{sec:Mod5_Sec1}
%%%%%%%%%%%%%%%%%%%%%%%%%%%%%%%%%%%%%%%%%%%%%%%%%%%%%%%%%%%%%%%%%%%%%%%
%

%%%%%%%%%%%%%%%%%%%%%%%%%%%%%%%%%%%%%%%%%%%%%%%%%%%%%%%%%%%%%%%%%%%%%%%
\section{Video naar afbeeldingen m.b.v. app VideoNaarAfbeeldingen}
\label{sec:Mod5_Sec2}
%%%%%%%%%%%%%%%%%%%%%%%%%%%%%%%%%%%%%%%%%%%%%%%%%%%%%%%%%%%%%%%%%%%%%%%
%

%%%%%%%%%%%%%%%%%%%%%%%%%%%%%%%%%%%%%%%%%%%%%%%%%%%%%%%%%%%%%%%%%%%%%%%
\section{Gemiddelde roodwaarde voor elke afbeelding berekenen}
\label{sec:Mod5_Sec3}
%%%%%%%%%%%%%%%%%%%%%%%%%%%%%%%%%%%%%%%%%%%%%%%%%%%%%%%%%%%%%%%%%%%%%%%
%

%%%%%%%%%%%%%%%%%%%%%%%%%%%%%%%%%%%%%%%%%%%%%%%%%%%%%%%%%%%%%%%%%%%%%%%
\section{Grafiek van de gemiddelde roodwaarde i.f.v. de tijd}
\label{sec:Mod5_Sec4}
%%%%%%%%%%%%%%%%%%%%%%%%%%%%%%%%%%%%%%%%%%%%%%%%%%%%%%%%%%%%%%%%%%%%%%%
%

%%%%%%%%%%%%%%%%%%%%%%%%%%%%%%%%%%%%%%%%%%%%%%%%%%%%%%%%%%%%%%%%%%%%%%%
\section{Automatische piekdetectie en berekening van de hartslag}
\label{sec:Mod5_Sec5}
%%%%%%%%%%%%%%%%%%%%%%%%%%%%%%%%%%%%%%%%%%%%%%%%%%%%%%%%%%%%%%%%%%%%%%%
%

%%%%%%%%%%%%%%%%%%%%%%%%%%%%%%%%%%%%%%%%%%%%%%%%%%%%%%%%%%%%%%%%%%%%%%%
\section{Extra: frequentie-analyse m.b.v. Fouriertransformatie}
\label{sec:Mod5_Sec6}
%%%%%%%%%%%%%%%%%%%%%%%%%%%%%%%%%%%%%%%%%%%%%%%%%%%%%%%%%%%%%%%%%%%%%%%
%


%%%%%%%%%%%%%%%%%%%%%%%%%%%%%%%%%%%%%%%%%%%%%%%%%%%%%%%%%%%%%%%%%%%%%%%
%%%%%%%%%%%%%%%%%%%%%%%%%%%%%%%%%%%%%%%%%%%%%%%%%%%%%%%%%%%%%%%%%%%%%%%
\deelmetoef{Module 6}{Experimenteren met de app}{Module 6. Experimenteren met de app}{Oplossingen module 6}{Oplossingen module 6}
%%%%%%%%%%%%%%%%%%%%%%%%%%%%%%%%%%%%%%%%%%%%%%%%%%%%%%%%%%%%%%%%%%%%%%%
%%%%%%%%%%%%%%%%%%%%%%%%%%%%%%%%%%%%%%%%%%%%%%%%%%%%%%%%%%%%%%%%%%%%%%%

\begin{samenvatting}
	Nu we een app hebben gemaakt waarmee we onze hartslag kunnen meten is het tijd om de app te evalueren. Net zoals bij elk ontwikkeld product zijn er wellicht verbeteringen mogelijk. In deze module gaan we met de app aan de slag en vergelijken we de app met alternatieven om je hartslag te meten, met als bedoeling eventuele fouten te achterhalen en verbeterpunten te ontdekken. Kritisch stilstaan bij je ontwerp en bij je product is noodzakelijk als je je product wilt blijven verbeteren.
\end{samenvatting}
%

%%%%%%%%%%%%%%%%%%%%%%%%%%%%%%%%%%%%%%%%%%%%%%%%%%%%%%%%%%%%%%%%%%%%%%%
\section{Invloed van sport en rust}
\label{sec:Mod6_Sec2}
%%%%%%%%%%%%%%%%%%%%%%%%%%%%%%%%%%%%%%%%%%%%%%%%%%%%%%%%%%%%%%%%%%%%%%%
%

Het is natuurlijk interessant om te weten wat je hartslag bij rust is, zoals besproken in Module \ref{sec:inl}, maar de meeste mensen gebruiken hartslagmeters om na te gaan wat je hartslag is bij inspanning.

\begin{opdracht}{Hartslag bij rust en na inspanning meten}
Meet je hartslag vlak voor en vlak na een inspanning. Zit je hartslag na inspanning binnen het aanbeloven hartslagbereik voor jouw leeftijd?

Verschilt je hartslag bij rust als je rechtop staat, rechtop zit of ligt? Waaraan zou dit kunnen liggen?
\end{opdracht}

%%%%%%%%%%%%%%%%%%%%%%%%%%%%%%%%%%%%%%%%%%%%%%%%%%%%%%%%%%%%%%%%%%%%%%%
\section{Vergelijking met bloeddrukmeter/hartslagmeter/smartwatch}
\label{sec:Mod6_Sec1}
%%%%%%%%%%%%%%%%%%%%%%%%%%%%%%%%%%%%%%%%%%%%%%%%%%%%%%%%%%%%%%%%%%%%%%%
%
Zoals we in Module \ref{sec:inl} besproken hebben, bestaan er alternatieven om je hartslag te meten. We spraken over bloeddrukmeters en hartslagmeters. Maar ook apps zoals de onze zijn al te verkrijgen in de Google Play Store.

\begin{opdracht}{Vergelijking met alternatieve hartslagmonitors}
Verzamel een aantal alternatieven voor onze hartslagmonitor-app. 

Som eerst op welke factoren je belangrijk vindt bij het gebruik van een hartslagmonitor. 

Ga na hoe goed de verschillende alternatieven naar jouw mening scoren voor elk van die factoren. Dit is een kwalitatieve evaluatie. Jij beoordeelt de kwaliteit van elk alternatief, op basis van jouw aanvoelen.

Misschien kan je sommige van die scores ook objectiever maken. Objectieve maatstaven zijn het tegengestelde van subjectieve maatstaven. Objectiever wil zeggen dat het oordeel minder afhankelijk is van jouw aanvoelen of van jouw persoonlijke mening, maar meer neutraal, zonder vooroordeel, kan vastgesteld worden. Vaak worden hiervoor kwantatieve maten gebruikt: dit zijn factoren die meetbaar zijn. 

Wellicht is de snelheid van de hartslagmeting voor jou belangrijk. Je subjectieve aanvoelen is wellicht dat de app die wij ontworpen hebben minder snel is dan een commercieel beschikbare hartslagmeter. (Dit is natuurlijk omdat het ook doenbaar moet zijn om de hartslagmonitor binnen een beperkt aantal lessen te maken, op basis van de kennis die jullie nu al hebben.) Maar als objectieve maat volstaat dit aanvoelen niet. Je zou met een chronometer kunnen nagaan hoelang de hartslagmeting duurt voor elk van de alternatieven. Dit is een objectieve maatstaf voor de snelheid van de hartslagmeting. 

Kan je ook voor de andere factoren die jij belangrijk vindt bepalen hoe je elke factor een objectieve score zou kunnen geven?
\end{opdracht}



%%%%%%%%%%%%%%%%%%%%%%%%%%%%%%%%%%%%%%%%%%%%%%%%%%%%%%%%%%%%%%%%%%%%%%%
\section{Evaluatie: wat kan beter?}
\label{sec:Mod6_Sec3}
%%%%%%%%%%%%%%%%%%%%%%%%%%%%%%%%%%%%%%%%%%%%%%%%%%%%%%%%%%%%%%%%%%%%%%%
%
\begin{opdracht}{Plus- en minputen en suggesties voor verbetering bepalen}
	Uit de vergelijking met de alternatieven heb je zelf kunnen ervaren hoe goed of slecht onze hartslagmonitor werkt. 

\begin{steroef}
	Op welke vlakken scoort de hartslagmonitor goed?
Wat kan beter? 
\end{steroef}

\begin{steroef}
	Als ontwerper moet je niet alleen bepalen wat de minpunten van je ontwerp zijn; je moet ook nadenken over hoe je het ontwerp kan verbeteren. Welke aanpassingen zou jij doen?
\end{steroef}

\begin{steroef}
	In een ideale wereld kan je uiteraard veel aanpassingen doen, maar als echte ontwerper ben je meestal meer beperkt. Welke beperkingen zie jij? 
\end{steroef}
\oplos{\begin{itemize}
		\item Budget
		\item Tijd
		\item Wiskundige achtergrond
		\item Beperkingen hardware: snelle vs trage smartphone, geheugen, wel of geen flash, \ldots
	\end{itemize}}

\begin{steroef}
	De beperkingen bepalen ook de haalbaarheid van de verbetering. Welke aanpassingen zijn haalbaar, welke zijn moeilijker te implementeren?
\end{steroef}

\begin{opmerking}
	Merk op dat voor dit probleem (net als voor veel andere problemen!) geen modeloplossing beschikbaar is. Vaak is een probleem \emph{open-ended}; dit wil zeggen dat niet op voorhand vastligt wat de beste oplossing voor het probleem is. Soms kan je door ervaring inschatten welke ontwerpen de beste zullen zijn, maar soms moet je het ook gewoon uitproberen en via evaluatie achterhalen wat goed en wat minder goed loopt en welke aanpassingen de kwaliteit van je producht ten goede be\"invloeden. Deze probleemoplossende vaardigheden zijn cruciaal in veel beroepen en vergen een zekere kritische houding en creativiteit.
\end{opmerking}
\end{opdracht}

%%%%%%%%%%%%%%%%%%%%%%%%%%%%%%%%%%%%%%%%%%%%%%%%%%%%%%%%%%%%%%%%%%%%%%%%%


%% Hier eindigt je eigenlijke LaTeX-file.

%% Van de volgende regel moet je afblijven.

\sluitfiles

%% Vervolgens kies je \'{e}\'{e}n van de onderstaande blokken, zonder iets aan deze blokken te veranderen:
%% ofwel produceer je achteraan de oplossingen van de oefeningen zonder ster; ofwel produceer je achteraan alle oplossingen.

%% Begin van het blok voor de oplossingen zonder ster %%

\deelzonderoef{}{Oplossingen van de oefeningen zonder $\ast$}{Oplossingen van de oefeningen zonder $\ast$}

\renewcommand{\sectionmark}[1]{\markright{\mijnoplosmark}}
\setcounter{secnumdepth}{0}
\titleformat{\section}[block]{}{}{0pt}{\zondernummer{#1}}

\input{inputs/oplossingen-zonder-ster.tmp}
%% Einde van het blok voor de oplossingen zonder ster %%

%% Begin van het blok voor alle oplossingen%%

\deelzonderoef{}{Oplossingen van alle oefeningen}{Oplossingen van alle oefeningen}

\renewcommand{\sectionmark}[1]{\markright{\mijnoplosmark}}
\setcounter{secnumdepth}{0}
\titleformat{\section}[block]{}{}{0pt}{\zondernummer{#1}}

\input{inputs/oplossingen-alle.tmp}
%% Einde van het blok voor alle oplossingen%%

%% In de volgende en laatste regel importeer je de achterflap. Je maakt een keuze tussen:
%% \includepdf{inputs/achterflap-leerling.pdf}
%% \includepdf{input/achterflap-leraar.pdf}
\includepdf{inputs/backcover.pdf}

\end{document} 